\documentclass[letterpaper,12pt]{article}
\begin{document}
%% \newcommand{\TBC}{\framebox{\textbf{TO BE COMPLETED}}}
\newcommand{\TBC}{}

\newtheorem{definition}{Definition}


\newcommand{\tet}{\mathtt{tet}}
\newcommand{\KC}{\mathtt{KC}}
\newcommand{\KV}{\mathtt{KV}}

\newcommand{\xb}{\mathtt{xb}}

\newcommand{\wq}{\mathtt{wq}}
\newcommand{\wk}{\mathtt{wk}}
\newcommand{\wv}{\mathtt{wv}}

\newcommand{\wrms}{\mathtt{wrms}}

\newcommand{\rmsnorm}{\mathtt{rmsnorm}}
\newcommand{\vecmatmul}{\mathtt{vecmatmul}}

\section{Background}
\subsection{Dimensions}
\label{dimensions}
The constants in Table~\ref{shapes} define the sizes of the various arrays
described later. These are from the {\tt Config struct}.
\begin{table}[hb]
  \begin{tabular}{|l|l|l|} \hline \hline 
    {\bf Math} & {\bf Code} & {\bf Explanation} \\ \hline 
    \(n_L\)  & {\tt n\_layers} & number of layers in the network  \\ \hline 
    \(n_P\) & XX &  \TBC \\ \hline 
    \(n_S\)  & {\tt seq\_len} & max sequence length \\ \hline 
    \(n_D\) & {\tt dim} &  transformer dimension \\ \hline 
     XXXX   & {\tt hidden\_dim} &  for ffn layers \\ \hline 
    \(n_V\) &  {\tt vocab\_size} &  vocabulary size \\ \hline 
    \(n_H\) & \verb+n_heads+ & number of query heads \\ \hline
    \({n_H}_{\mathtt{KV}}\)   & \verb+n_kv_heads+ & number of key/value heads \\ \hline
    \(n_{H_0}\) & \verb+n_heads * head_size+ & \TBC \\ \hline 
    \(n_{H_1}\) & \verb+n_kv_heads * head_size+ & \TBC \\ \hline 
    \hline 
\end{tabular}
    \caption{Dimensions}
    \label{shapes}
\end{table}

\subsection{Definitions}

\begin{definition}
  \label{1D_matrix}
  A 1D matrix is defined using a type {\tt T} and positive integer
\(n_X\), representing the dimensions/shape.
Let {\tt M} be a 1D matrix. Then, 
\begin{itemize}
\item \(\mathtt{M}[x]\) is a scalar of type {\tt T}
\end{itemize}

{\tt M} is stored as a linear array at address \(m_{\mathtt{M}}\)
  such that 
    \begin{itemize} 
\item Address of \(\mathtt{M}[x]  = m_{\mathtt{M}} + x\)
\end{itemize}
\end{definition}

\begin{definition}
  \label{2D_matrix}
  A 2D matrix is defined using a type {\tt T} and positive integers 
\(n_X, n_Y\), representing the dimensions/shape.
Let {\tt M} be a 2D matrix. Then, 
\begin{itemize}
\item \(\mathtt{M}[x]\) is a vector of length \(n_Y\)
\item \(\mathtt{M}[x, y]\) is a scalar of type {\tt T}
\end{itemize}

{\tt M} is stored as a linear array at address \(m_{\mathtt{M}}\)
  such that 
    \begin{itemize} 
\item Address of \(\mathtt{M}[x]  = 
  m_{\mathtt{M}} + (x \times n_Y)\)
\item Address of \(\mathtt{M}[x, y] = 
  m_{\mathtt{M}} + (x \times n_Y) + y\)
\end{itemize}
\end{definition}

\begin{definition}
  \label{3D_matrix}
  A 3D matrix is defined using a type {\tt T} and positive integers 
\(n_X, n_Y, n_Z\), representing the dimensions/shape.
Let {\tt M} be a 3D matrix. Then, 
\begin{itemize}
\item \(\mathtt{M}[x]\) is a 2D matrix of dimensions \(n_X, n_Y\)
\item \(\mathtt{M}[x, y]\) is a vector of length \(n_S\)
\item \(\mathtt{M}[x, y, z]\) is a scalar of type {\tt T}
\end{itemize}

{\tt M} is stored as a linear array at address \(m_{\mathtt{M}}\)
  such that 
    \begin{itemize} 
    \item Address of \(\mathtt{M}[x] = 
  m_{\mathtt{M}} + x \times (n_Y \times n_Z)\)
\item Address of \(\mathtt{M}[x, y]  = 
  m_{\mathtt{M}} + (x \times n_Y \times n_Z) + (y \times n_Z)\)
\item Address of \(\mathtt{M}[x, y, z] = 
  m_{\mathtt{M}} + (x \times n_X \times n_Y) + (y \times n_Z) + z\)
\end{itemize}
\end{definition}

\section{State}
\label{State}
The state is represented by the data structures listed in
Table~\ref{tbl_state_data}
\begin{table}
  \begin{tabular}{|l|l|l|l|l|l|} \hline \hline
    {\bf Name} & {\bf Shape} & {\bf Comment} & \(n_X\) & \(n_Y\) & \(n_Z\) \\ \hline
    {\tt x} & 1D array & activation at current timestamp &  \(n_D\) & & \\ \hline
    {\tt xb} & 1D array & like {\tt x} but inside residual branch & \(n_D\) & &
    \\ \hline
    {\tt xb2} & 1D array & like {\tt xb} & \(n_D\) & & \\ \hline
    {\tt KC} & 3D array & keys of kv-cache & \(n_L\) & \(n_P\) & \(n_S\)  \\
    \hline 
    {\tt KV} & 3D array & values of kv-cache & \(n_L\) & \(n_P\) & \(n_S\)  \\ \hline 
    \hline
\end{tabular}
  \caption{Data structures for state}
  \label{tbl_state_data}
\end{table}

\section{Weights}
\label{Weights}
The weights are represented by the data structures listed in
Table~\ref{tbl_weights_data}

\begin{table}
  \begin{tabular}{|l|l|l|l|l|l|} \hline \hline
    {\bf Name} & {\bf Shape} & {\bf Comment} & \(n_X\) & \(n_Y\) & \(n_Z\) \\ \hline
    {\tt tet} & 2D array & Lookup table maps token IDs to vectors  & \(n_V\) & \(n_D\) & \\ \hline 
    {\tt wk} & 3D array & key projections & \(n_L\) & \(n_D\) & \(n_{H_1}\) \\ \hline 
    {\tt wv} & 3D array & value projections & \(n_L\) & \(n_D\) & \(n_{H_1}\) \\ \hline 

    \hline
\end{tabular}
  \caption{Data structures for weights}
  \label{tbl_weights_data}
\end{table}

  
\section{Code}

Given 
\begin{enumerate}
\item an integer \(p\) where \(0 \leq p < n_P\)
\item an integer \(t\) where \(0 \leq t < n_V\)
  \end{enumerate}

%%--------------------------------------------
\begin{figure}
  \begin{tabbing} \hspace*{0.25in} \= \hspace*{0.25in} \=  \kill
\end{tabbing}
  \caption{Pseudo-code for \(\rmsnorm()\) function}
  \label{fn_rmsnorm}
\end{figure}
%%--------------------------------------------
\begin{figure}
  \begin{tabbing} \hspace*{0.25in} \= \hspace*{0.25in} \=  \kill
\end{tabbing}
  \caption{Pseudo-code for {\tt vecmatmul()} function}
  \label{fn_vecmatmul}
\end{figure}
%%--------------------------------------------
\begin{figure}
  \centering
\fbox{
  \begin{minipage}{8.5in}
  \begin{tabbing} 
    \hspace*{0.25in} \= \hspace*{0.25in} \=  
    \hspace*{0.25in} \= \hspace*{0.25in} \=  
    \kill
    {\bf function} {\tt foo} ({\bf integer} \(t\), {\bf integer} \(p\)) \+ \\
{\bf foreach} \(l \in \{0, 1, \ldots n_L-1\}\) {\bf do}  \+ \\ 

  %% copy the token embedding into x
  %% float* tet_ptr = mcr_2d_to_1d(w->token_embedding_table, token, ispc_dim);
  %% memcpy(x, tet_ptr, ((size_t)dim*sizeof(float)));
    \(x \leftarrow \tet[t]\) \\  
    % rmsnorm(s->xb, x, w_rms_att_l, dim);
    \(\xb \leftarrow \rmsnorm(x, \wrms[l])\) \\ 
    \(\mathtt{q}[l] \leftarrow \vecmatmul(\xb, \wq[l])\) \\ 
    \(\mathtt{k}[l] \leftarrow \vecmatmul(\xb, \wk[l])\) \\ 
    \(\mathtt{v}[l] \leftarrow \vecmatmul(\xb, \wv[l])\) \\ 
%%        matmul(s->q, s->xb, w->wq + l*dim*dim, dim, dim);
%%        matmul(s->k, s->xb, w->wk + l*dim*kv_dim, dim, kv_dim);
%%        matmul(s->v, s->xb, w->wv + l*dim*kv_dim, dim, kv_dim);
    \\ 
    {\bf endfor}  \- \\
    {\bf end} 
\end{tabbing}
  \end{minipage}
}
  \caption{Pseudo-code for {\tt forward()} function}
  \label{fn_forward}
\end{figure}
%      matmul(key_ptr, s->xb, w_k, dim, kv_dim);
%      matmul(val_ptr, s->xb, w_v, dim, kv_dim);

\section{Instructions to LLM}
For each of the functions listed in Table~\ref{list_of_fns}
\begin{enumerate}
  \item Write an optimized C function 
  \item Return a {\tt .c} file --- the definition
  \item Return a {\tt .h} file --- the declaration
\end{enumerate}
\begin{table}
\begin{tabular}{|l|l|} \hline \hline
  {\bf Function} & Section Implemented  \\ \hline \hline
  {\tt rmsnorm} & Section~\ref{fn_rmsnorm} \\ \hline
  {\tt vecmatmul} & Section~\ref{fn_vecmatmul} \\ \hline
  {\tt forward} & Section~\ref{fn_forward} \\ \hline
\end{tabular}
  \caption{Listing of functions}
  \label{list_of_fns}
\end{table}

\end{document}
